% C1 — Machine-verified rooted-tree majorants for polymer expansions
% with holes.  Charter: docs/C1-CHARTER.md (judges J-C1-1..5).
% REGIME: tricotomy everywhere — "exact" = Lean theorem compiled in the
% pinned core with clean axiom oracle; "certified" = committed
% interval-arithmetic transcript; "paper-level" = flagged prose.
% ZERO SLOTS as of the c1-v1.0 release (oracle + release facts inked
% from the recorded logs; see Reproducibility).
\documentclass[11pt]{article}
\usepackage[margin=1.1in]{geometry}
\usepackage{amsmath,amssymb,amsthm}
\usepackage{booktabs}
\usepackage[colorlinks=true,linkcolor=blue,urlcolor=blue,citecolor=blue]{hyperref}

\newtheorem{theorem}{Theorem}[section]
\newtheorem{lemma}[theorem]{Lemma}
\newtheorem{remark}[theorem]{Remark}
\newtheorem{definition}[theorem]{Definition}
\newtheorem{example}[theorem]{Example}

\newcommand{\lean}[1]{\texttt{#1}}
\newcommand{\ZL}{\mathbb{Z}/L}
\newcommand{\holes}{\mathcal{H}}
\newcommand{\skel}{\mathrm{sk}}
\newcommand{\Poly}{\Omega}

\title{Machine-checked rooted-tree majorants for\\
polymer expansions with holes}
\author{Lluis Eriksson\thanks{This manuscript and parts of the
formal development were produced with AI assistance under the
repository's audit protocol.  Every theorem below is machine-checked
in Lean~4 against a pinned Mathlib; the constant tables of
Section~\ref{sec:certified} are enclosed by committed interval
arithmetic; remaining numerics are labelled paper-level where they
occur.  Repository:
\url{https://github.com/lluiseriksson/THE-ERIKSSON-PROGRAMME}.}}
\date{July 2026}

\begin{document}
\maketitle

\begin{abstract}
We present a machine-checked quantitative toolkit for cluster
expansions of polymer systems with excluded regions
(\emph{holes}), in the discrete cube geometry of Balaban--Dimock
renormalization-group analyses.  Five Lean~4 theorems, checked
against a pinned Mathlib revision, provide:
(i) the identity $\sum_{T} \prod_v c_T(v)! = n!\,C_n$
for child factorials over spanning trees of the complete graph
$K_{n+1}$, with the rooted-tree majorant $4^n$ as corollary;
(ii) a marked-root leaf summation for the tree-graph majorant of an
Ursell-type expansion with holes, with the moment constant $M$ paid
once at the root and closed leaf ratio $4M^2$ per additional vertex,
together with its Catalan-sharpened form $M^{2n+1} C_n$ (a gain of
order $n^{3/2}$ in the $n$th coefficient); and
(iii) a target-preserving orderwise bound in which the target
union itself survives until the modified-metric exponential is
extracted.
A certified companion (interval arithmetic, 120-bit precision,
committed transcript with a committed reproduction witness) tabulates
the smallness gate and encloses every derived constant.  Non-vacuity
is machine-checked: a concrete hole family satisfying every
hypothesis is exhibited in Lean, and the two distinct hypothesis sets
among the polymer-facing theorems are both instantiated at it with a
strictly positive weight.  Each claim is labelled with its
verification layer: exact (Lean theorem), certified (interval
transcript), or paper-level.
\end{abstract}

\section{Introduction}

Cluster (polymer) expansions convert log-partition functions into
sums over connected geometric objects; their convergence proofs hinge
on combinatorial majorants for tree sums and on geometric smallness
gates \cite{KP86,FP07,Brydges86}.  When the expansion must respect
\emph{holes} --- excluded regions that the expansion sees as atomic
units, as in renormalization-group treatments of large-field regions
\cite{Balaban87,DimockII} --- the bookkeeping is notoriously
error-prone: root normalizations, child factorials, and
hole-separation hypotheses interact in ways that invite wrong
constants.\footnote{A wrong constant of exactly this kind occurred in
our own pre-registered prediction of the smallness-gate crossover;
the certified run caught it (Remark~\ref{rem:crossover}).}

This paper contributes, in machine-checked form:
\begin{enumerate}
\item an exact Catalan evaluation of the rooted child-factorial tree
  sum (Theorem~\ref{thm:catalanid}), of which the classical $4^n$
  tree-shape budget (Theorem~\ref{thm:tree}) is a corollary;
\item finite, volume-free bounds for the marked-root and
  fixed-target-union tree-graph majorants of an Ursell-type
  expansion with holes (Theorems~\ref{thm:markedroot},
  \ref{thm:markedrootcat}, \ref{thm:target}), with one smallness
  gate as the only inequality hypothesis and every constant an
  explicit function of $(d, \kappa_0)$;
\item a certified constant table (Section~\ref{sec:certified}) and a
  machine-checked non-vacuity witness (Section~\ref{sec:nonvacuity}).
\end{enumerate}
Statuses are kept separate throughout: \textbf{exact} means a Lean
theorem in the pinned core with axiom oracle
$[\lean{propext}, \lean{Classical.choice}, \lean{Quot.sound}]$;
\textbf{certified} means a committed interval-arithmetic transcript
with ball-plus-boolean output; \textbf{paper-level} means ordinary
mathematical prose not formalized here, and is always flagged.  The
formal development lives in a larger repository whose top-level
namespace (\lean{YangMills/}) reflects its original motivation in
lattice gauge theory; the present results are self-contained
combinatorial and expansion-theoretic tooling, and no reduction to
any continuum or gauge-theoretic problem is claimed or used.

\section{Setting and definitions}\label{sec:setting}

Throughout, $d, L \in \mathbb{N}$ with $L \ne 0$.

\begin{definition}[cubes and adjacency]
A \emph{cube} is a point of the discrete torus $(\ZL)^d$ (Lean:
\lean{Cube d L}).  Two cubes are \emph{adjacent} (\lean{cubeAdj})
iff they are distinct and each coordinate differs by $0$ or $\pm 1$;
this is the sup-metric neighbor graph, in which every cube has
$3^d - 1$ neighbors.  A cube set is \emph{connected} when it is
walk-connected in this graph.
\end{definition}

\begin{definition}[hole families and skeletons]\label{def:holes}
A \emph{hole family} $\holes$ (Lean: \lean{HoleFamily d L}) is a
finite set of finite cube sets --- the connected components of the
excluded region.  A cube set $X$ \emph{respects} $\holes$ iff every
hole $H \in \holes$ satisfies $H \subseteq X$ or $H \cap X =
\emptyset$: the expansion sees each hole as an atomic unit
\cite[eq.~(149)]{DimockII}.  The \emph{skeleton} $\skel(X)$ is the
set of cubes of $X$ lying in no hole.  The bounds below hypothesize:
pairwise disjointness of holes, no adjacency edge between distinct
holes, and non-emptiness of each hole.
\end{definition}

\begin{definition}[polymers]\label{def:polymers}
Given $\holes$ and an activity $z \colon \mathcal{P}((\ZL)^d)_{\rm
fin} \to \mathbb{C}$, a \emph{polymer} is a nonempty, connected,
hole-respecting cube set with nonempty skeleton; $\Poly$ denotes the
(finite) type of polymers (Lean: \lean{OmegaPolymerType}).  The
activity $z$ parameterizes the polymer type for compatibility with
the upstream expansion; no result below depends on it, and the
non-vacuity witness of Section~\ref{sec:nonvacuity} takes $z \equiv
0$.  Two
polymers are \emph{incompatible}, written $X \nsim X'$, iff their
skeletons intersect.  For a tuple $\mathbf{X} = (X_0, \dots, X_n) \in
\Poly^{n+1}$, its \emph{incompatibility graph} $G(\mathbf{X})$ has
vertex set $\{0,\dots,n\}$ and an edge $\{i,j\}$ iff $X_i \nsim X_j$.
\end{definition}

\begin{definition}[modified metric and hole weight]\label{def:metric}
The \emph{discrete modified metric} $\rho(X)$ (Lean:
\lean{discreteModifiedMetric}) is the Steiner length of the skeleton
inside $X$:
\[
\rho(X) \;=\; \min\bigl\{\, |S| - 1 \;:\; \skel(X) \subseteq S
\subseteq X,\ S \text{ connected} \,\bigr\},
\]
with $\rho(X) = 0$ if no such $S$ exists.  For a rate $\kappa \ge 0$
the \emph{hole weight} is
$W_\kappa(X) = e^{-\kappa\,(\rho(X)+1)}$
(Lean: \lean{appendixFHoleExpWeight}).
\end{definition}

\begin{definition}[the two tree-graph majorants]\label{def:sums}
Fix $\holes$, an activity $z$, a nonnegative weight $w \colon \Poly
\to \mathbb{R}_{\ge 0}$, and $n \in \mathbb{N}$.  Write
$\tau(\mathbf{X})$ for the number of spanning trees of
$G(\mathbf{X})$.  The \emph{marked-root sum} at a cube $r$ is
\[
S^{\#}_{w,r}(n) \;=\; \frac{1}{(n+1)!}
\sum_{\substack{\mathbf{X} \in \Poly^{n+1} \\ r \in \skel(X_0)}}
\tau(\mathbf{X}) \prod_{j=0}^{n} w(X_j)
\]
(Lean: \lean{appendixFHoleHsharpWeightedTreeMarkedRootSum}), and the
\emph{fixed-target sum} at a finite cube set $Y$ is
\[
T^{\#}_{w,Y}(n) \;=\; \frac{1}{(n+1)!}
\sum_{\substack{\mathbf{X} \in \Poly^{n+1} \\ \bigcup_j X_j = Y}}
\tau(\mathbf{X}) \prod_{j=0}^{n} w(X_j)
\]
(Lean: \lean{appendixFHoleHsharpWeightedTreeTerm}).
These are the standard tree-graph majorants of the order-$(n+1)$
coefficients of an Ursell (Mayer) expansion: by
tree-graph inequalities of Penrose type \cite{Penrose67,AR95}, Ursell
coefficients are dominated by spanning-tree counts of incompatibility
graphs, and it is the majorant objects $S^{\#}, T^{\#}$ --- not the
signed coefficients --- that the theorems below bound.  The passage
from signed coefficients to these majorants lives upstream in the
repository's Kotecký--Preiss/Mayer--Ursell development and is not
re-proved here.  (The word ``second'' in the Lean names refers to the
second cluster-expansion pass of the upstream development; nothing in
this paper depends on the first.)
\end{definition}

\begin{definition}[gate and derived constants]\label{def:constants}
The \emph{smallness gate} at dimension $d$ and rate $\kappa_0 > 0$ is
\begin{equation}\label{eq:gate}
G(d,\kappa_0) \;=\; (3^d)^2\, e^{-\kappa_0}\, 2^{3^d+1} \;<\; 1
\qquad (\text{Lean hypothesis } h_{Cq}),
\end{equation}
and the derived constants are
\begin{equation}\label{eq:constants}
\mathrm{RS} = (1-G)^{-1}, \qquad
M = \max\!\bigl(1,\ \kappa_0^{-1},\ \mathrm{RS}\bigr), \qquad
\Lambda = 4M^2
\end{equation}
(Lean: \lean{appendixFSecondUrsellMomentConstant},
\lean{appendixFSecondUrsellLeafConstant}).  All three are explicit
functions of $(d, \kappa_0)$ alone --- in particular
\emph{volume-free}: no constant below depends on the torus size $L$,
the target $Y$, or on $n$ beyond the displayed powers.  In
applications where the order-$n$ term additionally carries an
activity factor
$\varepsilon^n$, the resulting geometric series has ratio
$\Lambda\varepsilon$, which is below one exactly for $\varepsilon <
\varepsilon^* := \Lambda^{-1}$ (strict; nothing is claimed at
equality).  Status: the Lean endpoints prove the coefficient bounds
$M\Lambda^n$ and $M^{2n+1}C_n$ (exact); the enclosures of
$\Lambda^{-1}$ are certified; the passage from coefficient bound to a
convergent activity series is paper-level algebra conditional on the
stated $\varepsilon^n$ structure.
\end{definition}

\section{Main results}\label{sec:results}

All five statements carry stable, publication-facing names in module
\lean{MarkedRootedClosure}, file
\lean{PaperTheorems.lean}; each is a direct
application of a pinned upstream theorem --- no new axiom, no opaque
assumption.  Here $\mathcal{T}(K_{n+1})$ is the set of spanning trees
of the complete graph on the vertices $\{0, \dots, n\}$, each tree
rooted at $0$, and $c_T(v)$ is the number of children of $v$ in the
rooted tree $T$.

\begin{theorem}[exact Catalan value;
\lean{normalizedRootedChildFactorialTreeCatalanIdentity}]
\label{thm:catalanid}
For every $n \in \mathbb{N}$,
\[
\frac{n+1}{(n+1)!}
\sum_{T \in \mathcal{T}(K_{n+1})} \ \prod_{v} \bigl(c_T(v)\bigr)!
\;=\; C_n ,
\qquad\text{equivalently}\qquad
\sum_{T} \prod_v c_T(v)! \;=\; n!\,C_n ,
\]
where $C_n = \binom{2n}{n}/(n+1)$ is the $n$th Catalan number.
\end{theorem}

\begin{theorem}[rooted-tree majorant;
\lean{normalizedRootedChildFactorialTreeBound}]\label{thm:tree}
For every $n$, the left-hand side of Theorem~\ref{thm:catalanid} is
at most $4^n$.
\end{theorem}

Theorem~\ref{thm:tree} follows from Theorem~\ref{thm:catalanid} via
$C_n \le 4^n$; the classical budget overcounts by exactly the factor
$4^n/C_n$, which is $\sim \sqrt{\pi}\, n^{3/2}$.

\begin{theorem}[marked-root leaf summation;
\lean{markedRootLeafGeometricBound}]\label{thm:markedroot}
Let $\holes$ satisfy the separation hypotheses of
Definition~\ref{def:holes}, let $w \ge 0$ be dominated by the hole
weight at doubled rate, $w \le W_{2\kappa_0}$ pointwise, let $r$ be
any cube, and assume the gate \eqref{eq:gate}.  Then for every $n$,
\[
(n+1) \cdot S^{\#}_{w,r}(n) \;\le\; M \cdot \Lambda^{n}.
\]
\end{theorem}

\begin{theorem}[Catalan-sharpened marked-root bound;
\lean{markedRootCatalanBound}]\label{thm:markedrootcat}
Under exactly the hypotheses of Theorem~\ref{thm:markedroot},
\[
(n+1) \cdot S^{\#}_{w,r}(n) \;\le\; M^{2n+1} \cdot C_n .
\]
Since $M \Lambda^n = M^{2n+1} 4^n$, this improves
Theorem~\ref{thm:markedroot} by the factor $4^n/C_n$: the tree-shape
factor is paid at its exact Catalan value.
\end{theorem}

\begin{theorem}[target-preserving orderwise bound;
\lean{targetPreservingWeightedTreeBound}]\label{thm:target}
Let $\holes$ satisfy the separation hypotheses of
Definition~\ref{def:holes} and assume the gate \eqref{eq:gate}.  Let
$w, u \ge 0$ with $u \le W_{2\kappa_0}$ pointwise, suppose $w \le
W_{\mathrm{rate}} \cdot u$ pointwise for some rate $\ge 0$, and let
$Y$ be a finite cube set whose skeleton contains some cube $r$.  Then
\[
T^{\#}_{w,Y}(n) \;\le\; M \cdot W_{\mathrm{rate}}(Y)
\cdot \Lambda^{n}.
\]
Here ``orderwise'' means: for each fixed order $n$ separately.
\end{theorem}

\section{Proof architecture}\label{sec:architecture}

The formal proofs are the authority; this section gives the
mathematical mechanism at a level where a specialist can reconstruct
the logic.  Module pointers are to the repository.

\subsection{The Catalan identity}\label{sec:archcatalan}

The sum $\sum_T \prod_v c_T(v)!$ counts pairs $(T, \sigma)$ where $T$
is a spanning tree of $K_{n+1}$ rooted at $0$ and $\sigma$ assigns to
each vertex a linear order on its children --- i.e., it counts
\emph{labelled rooted plane trees} on the vertex set $\{0,\dots,n\}$
with root $0$.  Plane trees are rigid (a plane structure admits no
nontrivial automorphism fixing the root), so this count factors as
\[
\#\{\text{plane shapes with } n+1 \text{ nodes}\} \times
\#\{\text{labellings of the } n \text{ non-root nodes}\}
\;=\; C_n \cdot n! ,
\]
the first factor by the classical bijection between plane trees and
binary trees, formalized in \lean{YangMills/KP/PlaneTreeCatalan.lean}
through Mathlib's
\begin{center}\small
\lean{Tree.treesOfNumNodesEq\_card\_eq\_catalan};
\end{center}
\noindent the identity then normalizes to
Theorem~\ref{thm:catalanid}
(\lean{YangMills/KP/RootedCatalanExact.lean}).

\begin{example}
$n = 2$: $K_3$ has three spanning trees.  Rooted at $0$, the two
paths with $0$ as an endpoint contribute $\prod_v c_T(v)! = 1$ each,
and the path through $0$ --- the star at $0$ when rooted ---
contributes $2! = 2$; the sum is $4 = 2!\,C_2$.
\end{example}

\subsection{Leaf summation and the $2n+1$ moment budget}
\label{sec:archleaf}

Theorems~\ref{thm:markedroot} and \ref{thm:markedrootcat} are proved
by splitting the tuple sum by the (rooted) shape of a spanning tree
and peeling leaves.  Three mechanisms produce the constants:

\emph{One-pin polymer sums.}  The basic estimate controls the sum of
$W_{2\kappa_0}(X)$ over polymers $X$ whose skeleton contains a fixed
cube.  A polymer with $\rho(X) + 1 = m$ has a skeleton of at most $m$
cubes carried by a minimal connected Steiner set $S$; enumerating
rooted connected Steiner sets cube by cube costs $(3^d)^2$ per step
(degree-squared branching in the adjacency graph), and each cube of
$S$ carries a fiber entropy $2^{3^d+1}$ (the factor $2^{|S|}$ for
skeleton-subset choices times $2^{3^d|S|}$ for admissible hole
fillings over $S$).  The sum over polymers at Steiner length $m$ is
therefore dominated by $G^m$ with $G$ the gate \eqref{eq:gate}, and
the total by the geometric sum $\mathrm{RS} = (1-G)^{-1} \le M$.
The gate is the only inequality \emph{hypothesis}; the constants are
then explicit functions of $(d, \kappa_0)$, with the $\kappa_0^{-1}$
component of $M$ paid by the spare $\kappa_0$ of decay left over when
the doubled-rate weight $W_{2\kappa_0}$ surrenders only $W_{\kappa_0}$
to the geometric mechanism.

\emph{Moment paid once at the root.}  In the marked-root sum the
coordinate $X_0$ is pinned at $r$, so the root contributes one
one-pin sum: one factor $M$.  Each tree edge propagates a pin: a
child polymer must have skeleton intersecting its parent's, so
summing the child costs one one-pin sum ($M$) \emph{after} choosing
the pin inside the parent, whose cost is accounted to the parent's
own moment.  The bookkeeping closes with each of the $n$ non-root
vertices paying two factors ($M$ for its own one-pin sum, $M$ for the
pin it offers its children) and the root paying one: $M^{2n+1}$
total.  The formal count is the fixed-parent kernel lemma
\begin{center}\small
\lean{appendixFHoleHsharpWeightedTreeMarkedRootFixedParentKernelSum}\\
\lean{\_le\_prod\_childMomentBudget\_mul\_rootMoment}
\end{center}
and its consumers in
\lean{YangMills/RG/AppendixFSecondUrsellLeafSummation.lean}.

\emph{Shape sum.}  What remains is the sum over rooted tree shapes
with child multiplicities --- exactly the object of
Section~\ref{sec:archcatalan}.  Paying it at its Catalan value gives
Theorem~\ref{thm:markedrootcat}; paying it at the $4^n$ budget and
absorbing $4^n M^{2n} = \Lambda^n$ gives Theorem~\ref{thm:markedroot}.
The $(n+1)$ prefactor in both statements is the marked-coordinate
count: pinning coordinate $0$ rather than an arbitrary coordinate
costs a factor $n+1$ against the $(n+1)!$ normalization.

\subsection{Target preservation}\label{sec:archtarget}

In Theorem~\ref{thm:target} the tuple union is constrained to a fixed
target $Y$.  Extracting the metric exponential \emph{before} summing
would discard this constraint and with it the geometry of $Y$.
Instead, the hypothesis $w \le W_{\mathrm{rate}} \cdot u$ is applied
per polymer, and the extracted factors combine, for each contributing
tuple (connected in the incompatibility graph; disconnected tuples
contribute zero to $T^{\#}$), into at least the full-target factor
$W_{\mathrm{rate}}(Y)$: the superadditivity
$\rho(Y) + 1 \le \sum_i (\rho(X_i) + 1)$ is proved by stitching
minimal Steiner sets along a spanning tree of the incompatibility
graph (Lean:
\lean{omegaClusterUnion\_discreteModifiedMetric\_add\_one\_le}\allowbreak\lean{\_sum\_of\_spanningTree}).
Only after this factor is secured is the remainder $u$ routed through
the marked-root bound at a root $r \in \skel(Y)$.  The order of
operations is the content of the theorem's name and the reason the
constant is volume-free in $Y$.

\section{The certified companion}\label{sec:certified}

Script \lean{scripts/c1\_admissibility\_arb.py} (SHA-256
\begin{center}\small
\lean{f114fc4bcf2193c47603fb595b91b4e3ea62addc3a1bfebe0aa0c147c447f1f9}
\end{center}
\noindent \lean{python-flint} 0.9.0, precision 120 bits; transcript
committed alongside it, exit 0) encloses
$G$, $M$, $\Lambda$ (printed as \lean{Lf} in the transcript), and
$\varepsilon^*$ with midpoint--radius enclosures plus provable
booleans; the gate booleans are tri-state and provable-only
(\textsc{pass} only when $G < 1$ holds for every point of the
enclosing ball, \textsc{fail} only when $G \ge 1$ does,
\textsc{undecided} otherwise --- no \textsc{undecided} occurs in the
table).  The run was re-executed from scratch on 2026-07-11 and
reproduced the committed transcript byte-identically; the comparison
log (SHA-256 of both transcripts) is committed as
\lean{scripts/c1\_admissibility\_rerun\_20260711\_COMPARISON.txt}.
One header note: the script docstring's phrase ``certified activity
radius'' predates the status correction of
Definition~\ref{def:constants}, which governs.  Selected rows
(midpoints shown; radii in the transcript, all below $3 \times
10^{-6}$ relative):

\begin{center}
\begin{tabular}{@{}rrllll@{}}
\toprule
$d$ & $\kappa_0$ & gate $G$ & $M$ & $\Lambda$ & $\varepsilon^*$ \\
\midrule
2 & 11   & 1.3853059 & \multicolumn{3}{l}{\textsc{fail} (provably $\ge 1$)} \\
2 & 11.5 & 0.84023048 & 6.2590162 & 156.70114 & 0.0063815746 \\
2 & 12   & 0.50962555 & 2.0392580 & 16.634292 & 0.060116775 \\
2 & \textbf{16} & 0.0093341175 & 1.0094221 & 4.0757316 & \textbf{0.24535472} \\
2 & 24   & $3.13\times 10^{-6}$ & 1.0000031 & 4.0000251 & 0.24999843 \\
3 & 25   & 2.7177241 & \multicolumn{3}{l}{\textsc{fail} (provably $\ge 1$)} \\
3 & 26   & 0.99979481 & 4873.5020 & $9.50\times 10^{7}$ & $1.05\times 10^{-8}$ \\
3 & 28   & 0.13530751 & 1.1564805 & 5.3497886 & 0.18692327 \\
3 & \textbf{33} & 0.00091169486 & 1.0009125 & 4.0073035 & \textbf{0.24954436} \\
\bottomrule
\end{tabular}
\end{center}

The certified crossovers of the gate are
\[
\kappa^*(2) \in 11.32592096 \pm 2.72\times 10^{-10}, \qquad
\kappa^*(3) \in 25.99979479 \pm 2.32\times 10^{-9},
\]
and $\varepsilon^* \uparrow 1/4$ as $\kappa_0 \to \infty$ (the pure
tree-majorant limit $\Lambda \downarrow 4$).  The row $(d,\kappa_0) =
(3, 26)$ sits just above $\kappa^*(3)$ and displays the blow-up of
the derived constants as $G \to 1^-$: admissible, but with $M \approx
4.9 \times 10^{3}$ and $\varepsilon^* \approx 1.05 \times 10^{-8}$.

Scope of the certificates: the transcript's band booleans certify the
gate at the two registered band points $\kappa_0 = 16$ and
$\kappa_0 = 33$; the extension to all larger $\kappa_0$ is the
elementary monotonicity of $\kappa_0 \mapsto e^{-\kappa_0}$
(paper-level, one line).  Admissibility everywhere means the strict
gate $G < 1$, as in the Lean hypothesis.

\begin{remark}[a measured registration failure]\label{rem:crossover}
Before the certified run, we committed to the repository
(\lean{docs/C1-CHARTER.md}) a predicted crossover for the gate:
$15.72$ ($d{=}2$) and $32.60$ ($d{=}3$), from a hand-assembled
expression carrying $(3^d)^4$ where the gate \eqref{eq:gate} carries
$(3^d)^2$.  The certified run measured the discrepancy (transcript,
``crossover audit'': \textsc{registered prediction does not
re-derive}); the pre-committed admissibility bands $\kappa_0 \ge 16$
($d{=}2$) and $\kappa_0 \ge 33$ ($d{=}3$) remain true --- sufficient,
not sharp --- and only certified values appear in this paper.  We
record the failure as evidence for the correction protocol: constants
are computed, never assembled by hand.
\end{remark}

\section{What the Catalan sharpening buys}\label{sec:quantitative}

The gain of Theorem~\ref{thm:markedrootcat} over
Theorem~\ref{thm:markedroot} at order $n$ is the integer ratio
$4^n / C_n$ (computed in integer arithmetic, displayed rounded;
paper-level, independently recomputable):

\begin{center}
\begin{tabular}{@{}rrrr@{}}
\toprule
$n$ & $C_n$ & $4^n$ & $4^n/C_n$ \\
\midrule
1  & 1           & 4                & 4.00 \\
2  & 2           & 16               & 8.00 \\
5  & 42          & 1{,}024          & 24.38 \\
10 & 16{,}796    & 1{,}048{,}576    & 62.43 \\
20 & 6{,}564{,}120{,}420 & $\approx 1.10 \times 10^{12}$ & 167.50 \\
30 & $\approx 3.81 \times 10^{15}$ & $\approx 1.15 \times 10^{18}$ & 302.21 \\
\bottomrule
\end{tabular}
\end{center}

Being $\sim \sqrt{\pi}\,n^{3/2}$, the improvement is polynomial: it
does \emph{not} move the exponential threshold.  In an application
with per-vertex activity $\varepsilon$, both bounds give the same
radius requirement $4M^2\varepsilon < 1$ asymptotically; what changes
is every finite-order coefficient, hence every truncation-error
estimate at fixed order.  Concretely, at $(d, \kappa_0) = (2, 16)$
(certified row: $M = 1.0094221$, ball radius $\le 4 \times 10^{-8}$),
truncating so that order $n = 10$ is the last term retained, the
first omitted coefficient is smaller by the factor $4^{11}/C_{11}
\approx 71.4$ under the Catalan form --- almost two decimal digits of
truncation error at no analytic cost.  For the near-critical
certified row $(3, 26)$, where $M \approx 4.9 \times 10^3$ inflates
every coefficient through $M^{2n+1}$, the same polynomial factor
directly divides the order-$n$ term of any error budget.  These
consequences are paper-level arithmetic on top of the machine-checked
theorems.

\section{Non-vacuity, machine-checked}\label{sec:nonvacuity}

A characteristic failure mode of formalized bounds --- one this
programme has itself committed and repaired in earlier work --- is
hypotheses that no concrete geometry satisfies.  Module
\lean{MarkedRootedClosure/NonVacuity.lean} therefore exhibits --- as
instance terms, not existentials --- the hole family $\holes^{w}$ in
$d = 2$, $L = 8$ with two singleton holes at sup-distance $4$, and
proves by kernel decision (\lean{decide}) that the holes are
disjoint, edge-separated, and nonempty; the gate \eqref{eq:gate} at
$\kappa_0 = 16$ is proved in Lean via $e^{16} > 2.7^{16} > 82944$ ---
the Lean-exact counterpart of the certified row $G \approx 0.00933$.
The two distinct hypothesis sets among
Theorems~\ref{thm:markedroot}--\ref{thm:target} are then both
instantiated at $\holes^{w}$ with the \emph{strictly positive} weight
$W_{32} = e^{-32(\rho+1)}$ (no zero-weight degeneracy;
Theorem~\ref{thm:markedrootcat} shares its hypotheses verbatim with
Theorem~\ref{thm:markedroot}, so its satisfiability follows):
\begin{center}\small
\lean{nonvacuous\_markedRootLeafGeometricBound},\quad
\lean{nonvacuous\_targetPreservingWeightedTreeBound}.
\end{center}

What this witness proves: both hypothesis sets among
Theorems~\ref{thm:markedroot}--\ref{thm:target} are satisfiable by a
concrete geometry with positive weights, so the theorems have
non-trivial instances and the bounds are statements about nonzero
quantities.  What it does not prove: that the instantiated bounds are
sharp, or that the witness geometry models any physical large-field
configuration; it is a consistency certificate, not an application.

\section{Related work and novelty}\label{sec:related}

\emph{Cluster expansions.}  The abstract polymer framework and its
convergence criteria go back to Kotecký--Preiss \cite{KP86}, with the
modern sharp form of the criterion due to Fernández--Procacci
\cite{FP07}; pedagogical accounts include Brydges \cite{Brydges86}.
Tree-graph inequalities dominating Ursell coefficients by
spanning-tree counts originate with Penrose \cite{Penrose67}; the
family of tree and forest formulas is surveyed by
Abdesselam--Rivasseau \cite{AR95}.  The upstream repository
formalizes a Kotecký--Preiss criterion and a Mayer--Ursell expansion,
and the majorants $S^\#, T^\#$ of Definition~\ref{def:sums} are the
standard objects of this tradition.

\emph{Holes and large fields.}  The specific geometry ---
hole-respecting polymers, skeletons, and the modified (Steiner)
metric --- is that of renormalization-group analyses of lattice gauge
theory in the Balaban school \cite{Balaban87}, in the concrete form
of Dimock's expository series \cite{DimockI,DimockII,DimockIII}; the
hole-component families of Definition~\ref{def:holes} follow
\cite[\S 3.1.2]{DimockII}.  We are not aware of published
\emph{machine-checked} bounds for expansions in this geometry prior
to the present programme.

\emph{The Catalan identity.}  The count of labelled rooted plane
trees underlying Theorem~\ref{thm:catalanid} is classical
combinatorics at the level of \cite{Stanley15}, and we make no
novelty claim for the enumeration itself.  The identity in this
normalized child-factorial form, and its machine-checked proof, were
first announced in the author's own preprint \cite{Eriksson26} with
its companion repository; the present paper subsumes that
announcement into the hole-respecting expansion toolkit
(Theorem~\ref{thm:markedrootcat} and the certified constants).  We
are not aware of an \emph{independent} prior source for the identity
in exactly this form, for its use to replace the $4^n$ shape budget
inside an Ursell tree majorant with holes, or for a formalization of
either; readers aware of one are invited to point us to it.

\emph{What is new here.}  To our knowledge: (i) the formalized exact
evaluation (Theorem~\ref{thm:catalanid}) and its formalized corollary
chain into hole-respecting leaf summation
(Theorem~\ref{thm:markedrootcat}); (ii) the machine-checked
volume-free constant structure $M^{2n+1}$ with the single-gate
discipline \eqref{eq:gate}; (iii) the machine-checked non-vacuity
witness; and (iv) the certified-constants companion with committed
pre-registered bands and a documented registration failure
(Remark~\ref{rem:crossover}).  The mathematical ingredients are
classical; the contribution is the machine-checked integration.

\emph{Formalization context.}  The development builds on Mathlib
\cite{mathlib20}; the Catalan/plane-tree layer contributes upstream
API (a \lean{PlaneTree} companion prepared for separate submission to
Mathlib).

\section{Reproducibility}\label{sec:repro}

Toolchain: Lean \lean{leanprover/lean4:v4.29.0-rc6}; Mathlib pinned
to commit
\begin{center}\small
\lean{07642720480157414db592fa85b626dafb71355b}
\end{center}
\noindent (lakefile and manifest agree); \lean{python-flint} 0.9.0 on
CPython 3.12.6; \lean{tectonic} 0.15.0.

\medskip
\noindent Theorem-to-artifact map (statuses as in
Section~\ref{sec:results}; all Lean files in module
\lean{MarkedRootedClosure} unless noted):

\begin{center}
\footnotesize
\begin{tabular}{@{}llll@{}}
\toprule
Result & Lean name & File & Status \\
\midrule
Thm.~\ref{thm:catalanid} & \lean{...TreeCatalanIdentity} & \lean{PaperTheorems.lean} & exact \\
Thm.~\ref{thm:tree} & \lean{...TreeBound} & \lean{PaperTheorems.lean} & exact \\
Thm.~\ref{thm:markedroot} & \lean{markedRootLeafGeometricBound} & \lean{PaperTheorems.lean} & exact \\
Thm.~\ref{thm:markedrootcat} & \lean{markedRootCatalanBound} & \lean{PaperTheorems.lean} & exact \\
Thm.~\ref{thm:target} & \lean{targetPreservingWeightedTreeBound} & \lean{PaperTheorems.lean} & exact \\
\S\ref{sec:nonvacuity} witnesses & \lean{nonvacuous\_*} & \lean{NonVacuity.lean} & exact \\
\S\ref{sec:certified} table & --- & \lean{scripts/c1\_admissibility\_*} & certified \\
\S\ref{sec:quantitative} ratios & --- & integer arithmetic in text & paper-level \\
\S\ref{sec:quantitative} consequences & --- & --- & paper-level \\
\bottomrule
\end{tabular}
\end{center}

\medskip
\noindent From a clean clone of the repository at the release
revision:
\begin{verbatim}
lake exe cache get           # Mathlib oleans
lake build MarkedRootedClosure
lake env lean MarkedRootedClosure/Oracle.lean   # axiom oracle
python scripts/c1_admissibility_arb.py          # certified table
tectonic -X compile \
  papers/c1-rooted-tree-majorants/c1_rooted_tree_majorants.tex
\end{verbatim}
\noindent Release facts (recorded from the run logs):
\lean{lake build MarkedRootedClosure} completed with exit code $0$ in
\textbf{8244 jobs} at the Lean-integration commit
\begin{center}\small
\lean{4f95a5ad8154deef0ef38716073810072332fd37}
\end{center}
\noindent (release tag \lean{c1-v1.0} on the final manuscript
commit).  The axiom oracle
(\lean{lake env lean MarkedRootedClosure/Oracle.lean}; log committed
as \lean{scripts/c1\_oracle\_output.txt} at the same commit) reports,
for each of the seven theorems --- the five of
Section~\ref{sec:results} and the two non-vacuity instantiations of
Section~\ref{sec:nonvacuity} --- exactly
\begin{center}
\lean{[propext, Classical.choice, Quot.sound]},
\end{center}
\noindent the three standard classical axioms of Mathlib; no project
axiom, no \lean{sorry}.  The raw build transcripts (all three
recorded runs, including the two measured failures on the way to
green), the canonical-hash manifest distinguishing the integration
commit, the manuscript commit, and the tag, and a dual-hash (LF
blob / CRLF worktree) statement of the reproduction witness are
committed under \lean{scripts/c1\_lake\_build\_*} and
\lean{papers/c1-rooted-tree-majorants/RELEASE-MANIFEST.md} in the
documentation patch \lean{c1-v1.0.1}.

Pin-drift audit (measured): relative to the module's original
\lean{archive/UPSTREAM.lock} (upstream commit \lean{4e45246a\ldots}),
the blob of \lean{KP/RootedLeafSummation.lean} is byte-identical at
the build revision, while
\lean{RG/AppendixFSecondUrsellLeafSummation.lean} drifted
($+144$/$-15$ lines: the \lean{treeShape} refactor and the Catalan
route; the geometric theorem's statement is byte-unchanged).  The
Catalan chain (\lean{KP/RootedCatalanExact.lean},
\lean{KP/PlaneTreeCatalan.lean} and their dependencies) postdates the
lock and is therefore not audited against it; for those files, and
for everything else, the authority is the build itself: the theorems
are compiled against the release revision, not the lock.

\section{Conclusion}

Five machine-checked theorems give the exact combinatorial value and
the volume-free analytic envelope of rooted tree-graph majorants for
polymer expansions with holes, with certified constants and a
machine-checked consistency witness.  The limitations are stated
where they arise: the signed-coefficient-to-majorant passage is
upstream and the activity-series algebra is paper-level; the
Catalan gain is polynomial and moves no radius; the witness certifies
satisfiability, not sharpness.  Reasonable next uses are (a)
consuming Theorem~\ref{thm:markedrootcat} inside the upstream
programme wherever the $4^n$ budget is currently paid, and (b)
upstreaming the plane-tree API.

\begin{thebibliography}{99}

\bibitem{KP86} R.~Kotecký and D.~Preiss,
\emph{Cluster expansion for abstract polymer models},
Comm. Math. Phys. \textbf{103} (1986), 491--498.
doi:10.1007/BF01211762.

\bibitem{FP07} R.~Fernández and A.~Procacci,
\emph{Cluster expansion for abstract polymer models. New bounds from
an old approach}, Comm. Math. Phys. \textbf{274} (2007), 123--140.
doi:10.1007/s00220-007-0279-2, arXiv:math-ph/0605041.

\bibitem{Brydges86} D.~C.~Brydges,
\emph{A short course on cluster expansions}, in: Phénomènes critiques,
systèmes aléatoires, théories de jauge (Les Houches 1984),
North-Holland, Amsterdam, 1986, pp. 129--183.

\bibitem{Penrose67} O.~Penrose,
\emph{Convergence of fugacity expansions for classical systems}, in:
Statistical Mechanics: Foundations and Applications (T.~A.~Bak, ed.),
Benjamin, New York, 1967, pp. 101--109.

\bibitem{AR95} A.~Abdesselam and V.~Rivasseau,
\emph{Trees, forests and jungles: a botanical garden for cluster
expansions}, in: Constructive Physics (V.~Rivasseau, ed.), Lecture
Notes in Physics \textbf{446}, Springer, Berlin, 1995, pp. 7--36.
doi:10.1007/3-540-59190-7\_20.

\bibitem{Balaban87} T.~Bałaban,
\emph{Renormalization group approach to lattice gauge field theories.
I. Generation of effective actions in a small field approximation and
a coupling constant renormalization in four dimensions},
Comm. Math. Phys. \textbf{109} (1987), 249--301.

\bibitem{DimockI} J.~Dimock,
\emph{The renormalization group according to Balaban, I. Small
fields}, Rev. Math. Phys. \textbf{25} (2013), 1330010.
arXiv:1108.1335.

\bibitem{DimockII} J.~Dimock,
\emph{The renormalization group according to Balaban, II. Large
fields}, J. Math. Phys. \textbf{54} (2013), 092301.
arXiv:1212.5562.

\bibitem{DimockIII} J.~Dimock,
\emph{The renormalization group according to Balaban, III.
Convergence}, Ann. Henri Poincaré \textbf{15} (2014), 2133--2175.
arXiv:1304.0705.

\bibitem{Stanley15} R.~P.~Stanley,
\emph{Catalan Numbers}, Cambridge University Press, New York, 2015.

\bibitem{Eriksson26} L.~Eriksson,
\emph{A machine-verified bijective proof of the rooted
child-factorial Catalan identity over spanning trees of the complete
graph}, preprint, ai.viXra.org:2607.0001, July 2026.

\bibitem{mathlib20} The mathlib Community,
\emph{The Lean mathematical library}, in: Proc. 9th ACM SIGPLAN
Int. Conf. on Certified Programs and Proofs (CPP 2020), ACM, 2020,
pp. 367--381. doi:10.1145/3372885.3373824.

\end{thebibliography}

\end{document}
